\phantomsection
\chapter*{Ortak Veri ve Yazılım Uygulamaları}
\markboth{Ortak Veri ve Yazılım Uygulamaları}{Ortak Veri ve Yazılım Uygulamaları}
\addcontentsline{toc}{chapter}{Ortak Veri ve Yazılım Uygulamaları}
\label{sec:spss-guide}

Kitabın ortak literatür veri uygulamaları, \citet{cortez2008data} tarafından
paylaşılan \emph{Student Performance} veri setinin matematik dosyasını
kullanır. Veri, Portekiz'deki iki ortaöğretim okulunun öğrenci kayıtları ve
anketlerinden derlenmiştir. Yüklenen dosyada 395 satır ve 33 değişken vardır.
Böylece aynı araştırma bağlamı veri tanımından sonuç raporlamaya kadar izlenir;
her uygulama için farklı bir araştırmadan veri getirildiği iddia edilmez.

\begin{researchbox}
Bu veri iki okuldan gözlemsel kayıtlardır; Türkiye'deki öğrencilerin veya
Portekiz'deki bütün öğrencilerin olasılıklı örneklemi olarak sunulamaz.
Uygulamalardaki standart hatalar, güven aralıkları ve testler, bağımsız
benzer gözlemler modelinin öğretim amaçlı sonuçlarıdır. Okul ve sınıf içi
bağımlılık, seçilim ve karıştırıcı değişkenler gerçek bir araştırmada ayrıca
ele alınmalıdır. Dosyanın tamamını kullanmak, hedef evrenin tam sayımı
anlamına gelmez.
\end{researchbox}

\section*{Dosya düzeni ve çalıştırma}

Eşlikçi paketin \texttt{companion/spss} klasöründe
\texttt{csv/bXX.csv}, \texttt{excel/bXX.xlsx}, \texttt{syntax/bXX.sps},
\texttt{r/bXX.R} ve \texttt{python/bXX.py} bulunur; XX, 01--14'tür.
Excel dosyalarının \texttt{veri}, \texttt{sozluk} ve \texttt{kaynak}
sayfaları veri, sözlük ve kaynak bilgisini taşır. Üç dilin analiz kodları
ilgili bölümün basılı bloklarından üretilir. Uygulama kodları sabittir;
kapsamlı kitapta kod ile basılı bölüm numarası farklı olabilir.

Paketi bilgisayarınıza klasör yapısını koruyarak alın.
SPSS'te çalışma dizinini \texttt{companion} klasörünün bulunduğu proje kökü
olarak ayarlayın (\texttt{CD} komutunun dizinini kendi bilgisayarınıza göre
belirtin). \emph{File > Open > Syntax} ile ilgili \texttt{bXX.sps} dosyasını
açıp \emph{Run > All} seçin. B01, CSV açan \texttt{open-b01.sps}
dosyasını çağırır. B02/B03 de bu hazırlığı yeniden çalıştırır ve oluşan
\texttt{sav/b01.sav} dosyasını kullanır. B08'in \texttt{open-b08.sps}
dosyası Excel'in \texttt{veri} sayfasını açar; diğer güncel analizler
CSV'yi doğrudan açan bloklarla başlar. Etiketler ve ölçme düzeyleri
açma kodunda tanımlanır. Hazır SPSS veri/çıktı dosyası dağıtılmaz;
\texttt{.sav} dosyaları ilgili açma komutuyla SPSS'te oluşturulur.
Önceki analizden kalan filtre, ağırlık ve grup ayrımı kapatılır.
Her uygulamaya kendi komut dosyasını baştan çalıştırarak başlayın.
Komut sözdizimi için \citet{ibmspsssyntax} temel alınmıştır.

\textbf{Menü adımlarıyla çalışma.} Numaralı, açık zeminli kutular menüde
izlenecek işlemleri gösterir. Önce çalışma dizinini yukarıda belirtildiği
gibi ayarlayın. Bölümde ayrı veri açma bloğu varsa onu
\texttt{EXECUTE} dahil çalıştırın. B01--B03 için \texttt{open-b01.sps},
B08 için \texttt{open-b08.sps} kullanın. Bu hazırlık veriyi ve etiketleri
yükler, analizi başlatmaz. Ardından kutulu menü adımlarını izleyin.
Alternatif olarak \texttt{bXX.sps} dosyasının
tamamını çalıştırabilirsiniz; ayrıca menüden aynı analizi yapmanız gerekmez.
Özellikle seçim ve birleştirme sonrasında iki yolu karıştırmayın.

Kutularda İngilizce arayüz adları kullanılmıştır. \emph{Compare Means}
menüsü yeni sürümlerde \emph{Compare Means and Proportions} olarak
görünebilir. Sayısal girişlerde SPSS'in ondalık ayarına uygun ayırıcıyı
kullanın; verilen komutlarda ondalık nokta kullanılır. \emph{Continue}
alt penceredeki seçimi onaylar, \emph{OK} işlemi çalıştırır.
Menü açıklamaları IBM'in temel analiz ve kullanıcı kılavuzlarıyla birlikte
okunmalıdır \citep{ibmspssbase,ibmspsscore}. Çıktı kutularından sonraki kısa
raporların doğrulama kapsamı ilgili bölümde belirtilir. Paylaşılan
çıktılardan hazırlanmış özetler ile yalnız beklenen kontrol değerlerini
ayırt edin; bunlar burada yeni bir SPSS çalıştırıldığı anlamına gelmez.


\textbf{R ve Python ile çalışma.} Her uygulamanın R ve Python kodu,
\path{companion/spss/csv/bXX.csv} dosyasını okur. Bu dosya SPSS için
hazırlanan Excel kitabının \texttt{veri} sayfasıyla aynı kayıt ve sütunları
içerir; sözlük ve kaynak bilgisi Excel'in diğer sayfalarındadır. Çalışma
dizinini proje kökü olarak ayarlayıp kod bloklarını sırayla çalıştırın.
Eşlikçi \texttt{r/bXX.R} ve \texttt{python/bXX.py} dosyaları aynı blokları
birleştirir. R'deki \texttt{file.choose()} satırı yalnız yol sorunu için
alternatiftir; toplu çalıştırılan R dosyasına ikinci veri açma olarak eklenmez.
R örnekleri temel R ile \texttt{stats}, \texttt{utils} ve \texttt{graphics}
işlevlerini; Python örnekleri pandas, NumPy, SciPy ve gerektiğinde Matplotlib'i
kullanır. Python'daki test aralıkları için SciPy 1.11 veya sonrası gerekir.
İşlevlerin seçenekleri için resmi belgeler temel alınmıştır
\citep{rstatsmanual,scipystatsmanual}.

Ortak CSV'lerde eksik hücre yoktur; başlangıç kontrolleri eksik değer varsa
işlemi durdurur. Bu, farklı yazılımların sessizce farklı kayıtları analiz
etmesini önlemek içindir; yeni veride eksik değerlerin nasıl ele alınacağı
ayrıca kararlaştırılmalıdır. Sıfır notlar eksik sayılmaz. Benzetim verilerinde
hazır tekrarlar, örnekleme uygulamasında hazır seçim bayrakları kullanılır;
aynı tohumun farklı yazılımlarda aynı örneklemi üreteceği varsayılmaz.

Bu gerçek veri uygulamalarını, bölümdeki ayrı öğretim örneklerinden ayırın.
R ve SPSS'in burada yeniden çalıştırılmasıyla, kullanıcının paylaştığı
çıktıların karşılaştırılması aynı işlem değildir. Üç yazılımın
sonuçları aynı test türü, fark yönü, güven düzeyi ve düzeltme seçeneğiyle
karşılaştırılmalıdır. Hangi çıktının incelendiği bölüm içinde açıklanır.

\begin{mistakebox}
SPSS ve R bu ortamda çalıştırılmamıştır. İlgili bölümlerde kullanıcının
paylaştığı çıktılar Python hesaplarıyla karşılaştırılmıştır; kitap tabloları
bu çıktıların düzenlenmiş özetidir, yeni yazılım çıktısı değildir.
Excel dosyalarının XML/hücre eşitliği gerçek Excel/SPSS içe aktarma
denemesinin yerine geçmez. B01'de bildirilen Excel açma sorunu çözülmüş
sayılmaz; doğrulanan yol CSV'dir. B08'in Excel yolu da yerel SPSS'te
ayrıca denenmelidir. Menü ve tablo başlıkları sürüm/dile göre değişebilir.
\end{mistakebox}

\textbf{Değişkenleri doğru okuyun.} \texttt{id} yalnız dosya içi sıra
numarasıdır. \texttt{school}: 1=GP, 2=MS;
\texttt{sex}: kaynakta F ve M olarak verilen kategorilerin 1 ve 2 kodlarıdır.
\texttt{studytime} saat sayısı değil, sıralı süre kategorisidir.
\texttt{g1}, \texttt{g2}, \texttt{g3} sırasıyla birinci dönem, ikinci dönem ve
yıl sonu notlarıdır (0--20). \texttt{pass10}, öğretim için
$\texttt{g3}\geq10$ olduğunda 1, diğer durumda 0 olarak türetilmiştir;
kurumun bireysel öğrenci kararının yerine geçmez.

Kaynak dosyada boş hücre yoktur. Bununla birlikte bu durum bütün ölçümlerin
hatasız olduğunu göstermez. Yıl sonu notundaki 38 sıfır, eksik değer kodu
olarak değiştirilmemiştir. \texttt{absences} için kaynağın devamsızlık sayımı
korunur; gün veya saat birimi varsayılmaz. Benzetim uygulamalarında satırlar
öğrenci değil tekrar, örnekleme uygulamasında seçim bayrakları ve ağırlıklar
sonradan eklenen öğretim değişkenleridir.

\textbf{Kaynak ve yeniden üretim.} UCI kaydı veriyi CC BY 4.0 lisansıyla
sunar. Kaynak atfı, lisans ve dönüşümler her Excel dosyasına işlenmiştir.
\texttt{build\_package.py} veri paketini ve kontrol değerlerini yeniden
üretir; \texttt{results/checks.json} hesapları ve kaynak dosya özetlerini
saklar. Özgün 33 sütun \texttt{raw/student-mat.csv} içinde korunur.
Not: Daha sonra Portekizce ders dosyası da alınırsa dosyalar bağımsız iki
öğrenci grubu sayılmamalı, yapay \texttt{id} üzerinden eşleştirilmemelidir.